% P05.2 local rank-one graph orthogonality candidate, 25 September 2026.

Let $\mathbf X$ be a one-dimensional balanced rational polyhedral space,
and fix a finite rational polyhedral structure $\Pi$ carrying its balancing
cycle.  For a vertex $v\in\Pi^{(0)}$, write $E(v)$ for the incident edges,
$u_{e/v}$ for the primitive outgoing integral direction of $e$ at $v$, and
$w_e>0$ for the balancing weight of $e$.  Put
\[
 K_v:=\left\{\lambda\in\mathbf R_{\geq0}^{E(v)}:\
             \sum_{e\in E(v)}\lambda_eu_{e/v}=0\right\}.
\]
The local weight vector $w_v=(w_e)_{e\in E(v)}$ belongs to the relative
interior of $K_v$.

\begin{theorem}
\label{thm:frontier-polyhedral-local-rank-one-graph}
Assume that $\dim K_v=1$ for every vertex $v\in\Pi^{(0)}$.  Let
$\gamma\in\mathrm{PA}_{\mathbf R}(\mathbf X)\cap
\mathrm{PC}(\mathbf X)$ be positive at infinity.  Then
$(\mathbf X,\gamma)$ has real-polyhedral orthogonality: for every
$u\in\mathrm{PA}_{\mathbf R,b}(\mathbf X)$,
\[
 \int_{\mathbf X}
 \bigl(P_\gamma(u)-(\gamma+u)\bigr)
 \MA_{\mathrm{poly},\mathbf R}(P_\gamma(u))=0.
\]
The conclusion allows bounded edges and cycles in the underlying graph.
It depends only on the local cones $K_v$ and is unchanged by subdividing
$\Pi$.
\end{theorem}

\begin{corollary}
\label{cor:frontier-polyhedral-trivalent-graph}
Every one-dimensional balanced rational polyhedral space admitting a finite
structure whose vertices have valence at most three satisfies the conclusion
of \Cref{thm:frontier-polyhedral-local-rank-one-graph} for every
real-piecewise-affine polyhedrally convex reference positive at infinity.
In particular, lattice saturation and the absence of bounded cycles are not
required.
\end{corollary}

The theorem is a sufficient criterion, not a converse.  At a vertex with
$\dim K_v>1$, the one-vertex construction gives a local noncontact mass
mechanism, but extending that germ through all bounded edges while retaining
global polyhedral convexity and the prescribed growth is a separate
realizability problem.
